% Capítulo: Cumplimiento de Objetivos Específicos (para \input en la tesis)
% Nota: Este archivo no incluye preámbulo ni \begin{document}. Inserte con % Capítulo: Cumplimiento de Objetivos Específicos (para \input en la tesis)
% Nota: Este archivo no incluye preámbulo ni \begin{document}. Inserte con % Capítulo: Cumplimiento de Objetivos Específicos (para \input en la tesis)
% Nota: Este archivo no incluye preámbulo ni \begin{document}. Inserte con % Capítulo: Cumplimiento de Objetivos Específicos (para \input en la tesis)
% Nota: Este archivo no incluye preámbulo ni \begin{document}. Inserte con \input{docs/achievements/objetivos_especificos.tex}

\section*{Cumplimiento de Objetivos Específicos}
\noindent\textbf{Fecha:} 12 de noviembre de 2025\\

Este capítulo documenta, con formato de tesis, cómo se cumplieron los cuatro objetivos específicos del proyecto. Para cada objetivo se presenta: enunciado, enfoque, implementación, evidencia, métricas, incidencias y estado.

\subsection*{Lista de objetivos}
\begin{enumerate}
  \item Modelar la arquitectura de un prototipo de SMA capaz de orquestar y gestionar fases clave del pipeline de ML supervisado de clasificación (ingesta, preprocesamiento, preparación de datos) en entornos simulados.
  \item Implementar el prototipo de arquitectura SMA modelado, integrando al menos dos algoritmos representativos (Random Forest y SVM) para su ejecución durante la fase de entrenamiento, usando JADE como entorno multiagente.
  \item Evaluar la eficiencia computacional y la sostenibilidad energética del prototipo SMA frente a una arquitectura centralizada, durante el entrenamiento, mediante la medición de tiempo, throughput, uso de CPU y memoria, y energía normalizada (J/MB) en escenarios simulados, controlados y replicables.
  \item Analizar e interpretar los resultados para validar la eficiencia y sostenibilidad del SMA; sintetizar implicaciones técnicas, éticas y operativas; identificar desafíos, limitaciones y oportunidades de mejora.
\end{enumerate}

\subsection*{OE1 — Modelado de la arquitectura SMA}
\textbf{Enunciado.} Modelar la arquitectura de un prototipo de SMA para orquestar el pipeline de ML supervisado en entornos simulados.\\
\textbf{Enfoque.} Arquitectura JADE + Backend (FastAPI) con exposición de métricas; simulación de etapas PREPROCESS $\rightarrow$ TRAIN $\rightarrow$ EVAL con eventos JSONL.\\
\textbf{Implementación.}
\begin{itemize}
  \item Documentación de arquitectura: \texttt{docs/architecture/architecture.md}, índice \texttt{docs/architecture/README.md}.
  \item Puente con JADE: \texttt{docs/architecture/bridge\_jade.md}; protocolo de mensajes/eventos: \texttt{docs/architecture/message-protocol.md}.
  \item Scripts: \texttt{scripts/run\_jade\_experiment.sh}, \texttt{scripts/start\_stack.sh}.
\end{itemize}
\textbf{Evidencia.} Artefactos de orquestación en \texttt{data/results/\_jade\_logs/*}; métricas de progreso en \texttt{backend/api/main.py} (\texttt{/metrics}).\\
\textbf{Métricas.} \texttt{experiment\_last\_id}, \texttt{experiment\_current\_stage\_idx}, \texttt{experiment\_current\_stage\_elapsed\_seconds}.\\
\textbf{Incidencias.} Síntesis de artefactos JADE cuando faltan; normalización de documentos movidos a \texttt{docs/architecture/*}.\\
\textbf{Estado.} Cumplido.

\subsection*{OE2 — Implementación con RF y SVM (JADE)}
\textbf{Enunciado.} Implementar el prototipo SMA integrando Random Forest y SVM durante el entrenamiento con JADE.\\
\textbf{Enfoque.} Dataset reproducible y simulación de PREPROCESS/ TRAIN\_RF/ EVAL\_RF/ TRAIN\_SVM/ EVAL\_SVM, registrando eventos JSONL; JADE opera y/o se sintetizan artefactos.\\
\textbf{Implementación.}
\begin{itemize}
  \item Generador: \texttt{python-analysis/generate\_dataset.py} $\rightarrow$ \texttt{data/raw/synthetic\_classification.csv} y metadatos (SHA-256).
  \item Simulación: \texttt{python-analysis/simulate\_experiment.py}; reportes: \texttt{python-analysis/report\_run.py}.
  \item Plataforma JADE y lanzador: \texttt{agents/jade-platform/*}, \texttt{scripts/run\_jade\_experiment.sh}.
\end{itemize}
\textbf{Evidencia.} \texttt{data/results/\textless exp\_id\textgreater/} con \texttt{meta.json}, \texttt{logs/events.jsonl}, \texttt{report/report.html}.\\
\textbf{Métricas.} \texttt{train\_duration\_seconds\_last\_by\_algo}, \texttt{train\_cpu\_avg\_percent\_last\_by\_algo}, \texttt{train\_mem\_peak\_mb\_last\_by\_algo}, \texttt{train\_records\_per\_s\_last\_by\_algo}.\\
\textbf{Incidencias.} Variantes históricas \texttt{energy\_j\_*} normalizadas en \texttt{aggregate\_metrics.py}; CLASSPATH/JADE documentado en el script.\\
\textbf{Estado.} Cumplido.

\subsection*{OE3 — Evaluación de eficiencia y energía (SMA vs. centralizado)}
\textbf{Enunciado.} Evaluar eficiencia y sostenibilidad del SMA frente a una arquitectura centralizada durante el entrenamiento, con métricas (tiempo, throughput, CPU, memoria, energía J/MB) en escenarios simulados y replicables.\\
\textbf{Enfoque.} Comparación control (invocación directa) vs. tratamiento (orquestación) como proxy de arquitectura centralizada vs. SMA; dataset congelado y semillas controladas.\\
\textbf{Implementación.}
\begin{itemize}
  \item Ejecuciones reproducibles: \texttt{scripts/monitoring\_smoke.sh} y \texttt{scripts/run\_experiment.sh}.
  \item Agregación: \texttt{python-analysis/aggregate\_metrics.py} $\rightarrow$ \texttt{data/results/aggregate\_metrics.csv} con duración, CPU, memoria, throughput, \texttt{avg\_watts}, \texttt{watts\_per\_mb}, \texttt{watts\_per\_record} y Joules derivados (\\$J = W \times s\\$).
  \item Observabilidad: Prometheus + Grafana (\texttt{infra/compose/*}); dashboards versionados; dashboard mínimo en \texttt{docs/grafana/thesis\_minimal.json}.
\end{itemize}
\textbf{Evidencia.} CSV por grupo/algoritmo; paneles de performance/energía; estado del backend en \texttt{data/results/\_server\_state/*}.\\
\textbf{Métricas.} Eficiencia: \texttt{train\_duration\_seconds}, \texttt{train\_records\_per\_s\_last}, \texttt{train\_cpu\_avg\_percent\_last}, \texttt{train\_mem\_peak\_mb\_last}. Energía: \texttt{train\_avg\_watts\_last}, \texttt{watts\_per\_mb}, \texttt{watts\_per\_record}, derivación a \texttt{joules\_per\_mb}.\\
\textbf{Incidencias.} RAPL inaccesible $\rightarrow$ proxy de potencia \texttt{cpu\_power\_w\_proxy} (calibrable con \texttt{CPU\_POWER\_W}); targets Prometheus en modo host y puertos ocupados resueltos por scripts.\\
\textbf{Limitaciones.} La arquitectura “centralizada” se modela vía control vs. tratamiento; \texttt{joules\_per\_mb} se computa a partir de campos disponibles (posible añadir al CSV).\\
\textbf{Estado.} Cumplido (en alcance de prototipo).

\subsection*{OE4 — Análisis e interpretación de resultados}
\textbf{Enunciado.} Analizar e interpretar resultados para validar la eficiencia y sostenibilidad del SMA; sintetizar implicaciones técnicas, éticas y operativas; identificar desafíos/limitaciones/oportunidades.\\
\textbf{Enfoque.} Uso de CSV consolidado + reportes por ejecución + paneles de Grafana; discusión de amenazas a la validez y ética energética.\\
\textbf{Implementación.}
\begin{itemize}
  \item Consolidación: \texttt{aggregate\_metrics.py} y \texttt{data/results/aggregate\_metrics.csv}; estadísticas: \texttt{stats\_analysis.py}.
  \item Reportes por ejecución: \texttt{python-analysis/report\_run.py}.
  \item Documentación de análisis y validez: \texttt{docs/achievements/cumplimiento\_objetivo\_general.md} (apéndices), \texttt{docs/monitoring\_spec.md}, \texttt{docs/energy\_plan.md}.
\end{itemize}
\textbf{Evidencia.} Comparaciones control/tratamiento por algoritmo; discusión de incidencias reales y amenazas a la validez.\\
\textbf{Síntesis.} Técnicas (overhead vs. eficiencia; calibración de potencia), éticas (transparencia energética), operativas (scripts reproducibles, dashboards versionados).\\
\textbf{Desafíos.} Integrar RAPL por ventana en CSV; exponer \texttt{joules\_per\_mb}; extender a datasets reales y cargas intensivas.\\
\textbf{Estado.} Cumplido.

\subsection*{Conclusión}
Los cuatro objetivos específicos se encuentran cumplidos en el alcance de prototipo, con trazabilidad a código, scripts y artefactos, y con análisis cuantitativo y cualitativo suficiente para sustentar la evaluación de eficiencia y sostenibilidad.

\subsection*{Referencias cruzadas}
\begin{itemize}
  \item Arquitectura: \texttt{docs/architecture/README.md}, \texttt{architecture.md}, \texttt{bridge\_jade.md}, \texttt{message-protocol.md}.
  \item Observabilidad/energía: \texttt{docs/monitoring\_spec.md}, \texttt{docs/energy\_plan.md}, dashboards en \texttt{infra/compose/grafana/dashboards/}, \texttt{docs/grafana/README.md}.
  \item Cumplimiento objetivo general: \texttt{docs/achievements/cumplimiento\_objetivo\_general.md}.
  \item Pipelines/simulación/análisis: \texttt{python-analysis/*.py}, \texttt{scripts/*.sh}.
\end{itemize}


\section*{Cumplimiento de Objetivos Específicos}
\noindent\textbf{Fecha:} 12 de noviembre de 2025\\

Este capítulo documenta, con formato de tesis, cómo se cumplieron los cuatro objetivos específicos del proyecto. Para cada objetivo se presenta: enunciado, enfoque, implementación, evidencia, métricas, incidencias y estado.

\subsection*{Lista de objetivos}
\begin{enumerate}
  \item Modelar la arquitectura de un prototipo de SMA capaz de orquestar y gestionar fases clave del pipeline de ML supervisado de clasificación (ingesta, preprocesamiento, preparación de datos) en entornos simulados.
  \item Implementar el prototipo de arquitectura SMA modelado, integrando al menos dos algoritmos representativos (Random Forest y SVM) para su ejecución durante la fase de entrenamiento, usando JADE como entorno multiagente.
  \item Evaluar la eficiencia computacional y la sostenibilidad energética del prototipo SMA frente a una arquitectura centralizada, durante el entrenamiento, mediante la medición de tiempo, throughput, uso de CPU y memoria, y energía normalizada (J/MB) en escenarios simulados, controlados y replicables.
  \item Analizar e interpretar los resultados para validar la eficiencia y sostenibilidad del SMA; sintetizar implicaciones técnicas, éticas y operativas; identificar desafíos, limitaciones y oportunidades de mejora.
\end{enumerate}

\subsection*{OE1 — Modelado de la arquitectura SMA}
\textbf{Enunciado.} Modelar la arquitectura de un prototipo de SMA para orquestar el pipeline de ML supervisado en entornos simulados.\\
\textbf{Enfoque.} Arquitectura JADE + Backend (FastAPI) con exposición de métricas; simulación de etapas PREPROCESS $\rightarrow$ TRAIN $\rightarrow$ EVAL con eventos JSONL.\\
\textbf{Implementación.}
\begin{itemize}
  \item Documentación de arquitectura: \texttt{docs/architecture/architecture.md}, índice \texttt{docs/architecture/README.md}.
  \item Puente con JADE: \texttt{docs/architecture/bridge\_jade.md}; protocolo de mensajes/eventos: \texttt{docs/architecture/message-protocol.md}.
  \item Scripts: \texttt{scripts/run\_jade\_experiment.sh}, \texttt{scripts/start\_stack.sh}.
\end{itemize}
\textbf{Evidencia.} Artefactos de orquestación en \texttt{data/results/\_jade\_logs/*}; métricas de progreso en \texttt{backend/api/main.py} (\texttt{/metrics}).\\
\textbf{Métricas.} \texttt{experiment\_last\_id}, \texttt{experiment\_current\_stage\_idx}, \texttt{experiment\_current\_stage\_elapsed\_seconds}.\\
\textbf{Incidencias.} Síntesis de artefactos JADE cuando faltan; normalización de documentos movidos a \texttt{docs/architecture/*}.\\
\textbf{Estado.} Cumplido.

\subsection*{OE2 — Implementación con RF y SVM (JADE)}
\textbf{Enunciado.} Implementar el prototipo SMA integrando Random Forest y SVM durante el entrenamiento con JADE.\\
\textbf{Enfoque.} Dataset reproducible y simulación de PREPROCESS/ TRAIN\_RF/ EVAL\_RF/ TRAIN\_SVM/ EVAL\_SVM, registrando eventos JSONL; JADE opera y/o se sintetizan artefactos.\\
\textbf{Implementación.}
\begin{itemize}
  \item Generador: \texttt{python-analysis/generate\_dataset.py} $\rightarrow$ \texttt{data/raw/synthetic\_classification.csv} y metadatos (SHA-256).
  \item Simulación: \texttt{python-analysis/simulate\_experiment.py}; reportes: \texttt{python-analysis/report\_run.py}.
  \item Plataforma JADE y lanzador: \texttt{agents/jade-platform/*}, \texttt{scripts/run\_jade\_experiment.sh}.
\end{itemize}
\textbf{Evidencia.} \texttt{data/results/\textless exp\_id\textgreater/} con \texttt{meta.json}, \texttt{logs/events.jsonl}, \texttt{report/report.html}.\\
\textbf{Métricas.} \texttt{train\_duration\_seconds\_last\_by\_algo}, \texttt{train\_cpu\_avg\_percent\_last\_by\_algo}, \texttt{train\_mem\_peak\_mb\_last\_by\_algo}, \texttt{train\_records\_per\_s\_last\_by\_algo}.\\
\textbf{Incidencias.} Variantes históricas \texttt{energy\_j\_*} normalizadas en \texttt{aggregate\_metrics.py}; CLASSPATH/JADE documentado en el script.\\
\textbf{Estado.} Cumplido.

\subsection*{OE3 — Evaluación de eficiencia y energía (SMA vs. centralizado)}
\textbf{Enunciado.} Evaluar eficiencia y sostenibilidad del SMA frente a una arquitectura centralizada durante el entrenamiento, con métricas (tiempo, throughput, CPU, memoria, energía J/MB) en escenarios simulados y replicables.\\
\textbf{Enfoque.} Comparación control (invocación directa) vs. tratamiento (orquestación) como proxy de arquitectura centralizada vs. SMA; dataset congelado y semillas controladas.\\
\textbf{Implementación.}
\begin{itemize}
  \item Ejecuciones reproducibles: \texttt{scripts/monitoring\_smoke.sh} y \texttt{scripts/run\_experiment.sh}.
  \item Agregación: \texttt{python-analysis/aggregate\_metrics.py} $\rightarrow$ \texttt{data/results/aggregate\_metrics.csv} con duración, CPU, memoria, throughput, \texttt{avg\_watts}, \texttt{watts\_per\_mb}, \texttt{watts\_per\_record} y Joules derivados (\\$J = W \times s\\$).
  \item Observabilidad: Prometheus + Grafana (\texttt{infra/compose/*}); dashboards versionados; dashboard mínimo en \texttt{docs/grafana/thesis\_minimal.json}.
\end{itemize}
\textbf{Evidencia.} CSV por grupo/algoritmo; paneles de performance/energía; estado del backend en \texttt{data/results/\_server\_state/*}.\\
\textbf{Métricas.} Eficiencia: \texttt{train\_duration\_seconds}, \texttt{train\_records\_per\_s\_last}, \texttt{train\_cpu\_avg\_percent\_last}, \texttt{train\_mem\_peak\_mb\_last}. Energía: \texttt{train\_avg\_watts\_last}, \texttt{watts\_per\_mb}, \texttt{watts\_per\_record}, derivación a \texttt{joules\_per\_mb}.\\
\textbf{Incidencias.} RAPL inaccesible $\rightarrow$ proxy de potencia \texttt{cpu\_power\_w\_proxy} (calibrable con \texttt{CPU\_POWER\_W}); targets Prometheus en modo host y puertos ocupados resueltos por scripts.\\
\textbf{Limitaciones.} La arquitectura “centralizada” se modela vía control vs. tratamiento; \texttt{joules\_per\_mb} se computa a partir de campos disponibles (posible añadir al CSV).\\
\textbf{Estado.} Cumplido (en alcance de prototipo).

\subsection*{OE4 — Análisis e interpretación de resultados}
\textbf{Enunciado.} Analizar e interpretar resultados para validar la eficiencia y sostenibilidad del SMA; sintetizar implicaciones técnicas, éticas y operativas; identificar desafíos/limitaciones/oportunidades.\\
\textbf{Enfoque.} Uso de CSV consolidado + reportes por ejecución + paneles de Grafana; discusión de amenazas a la validez y ética energética.\\
\textbf{Implementación.}
\begin{itemize}
  \item Consolidación: \texttt{aggregate\_metrics.py} y \texttt{data/results/aggregate\_metrics.csv}; estadísticas: \texttt{stats\_analysis.py}.
  \item Reportes por ejecución: \texttt{python-analysis/report\_run.py}.
  \item Documentación de análisis y validez: \texttt{docs/achievements/cumplimiento\_objetivo\_general.md} (apéndices), \texttt{docs/monitoring\_spec.md}, \texttt{docs/energy\_plan.md}.
\end{itemize}
\textbf{Evidencia.} Comparaciones control/tratamiento por algoritmo; discusión de incidencias reales y amenazas a la validez.\\
\textbf{Síntesis.} Técnicas (overhead vs. eficiencia; calibración de potencia), éticas (transparencia energética), operativas (scripts reproducibles, dashboards versionados).\\
\textbf{Desafíos.} Integrar RAPL por ventana en CSV; exponer \texttt{joules\_per\_mb}; extender a datasets reales y cargas intensivas.\\
\textbf{Estado.} Cumplido.

\subsection*{Conclusión}
Los cuatro objetivos específicos se encuentran cumplidos en el alcance de prototipo, con trazabilidad a código, scripts y artefactos, y con análisis cuantitativo y cualitativo suficiente para sustentar la evaluación de eficiencia y sostenibilidad.

\subsection*{Referencias cruzadas}
\begin{itemize}
  \item Arquitectura: \texttt{docs/architecture/README.md}, \texttt{architecture.md}, \texttt{bridge\_jade.md}, \texttt{message-protocol.md}.
  \item Observabilidad/energía: \texttt{docs/monitoring\_spec.md}, \texttt{docs/energy\_plan.md}, dashboards en \texttt{infra/compose/grafana/dashboards/}, \texttt{docs/grafana/README.md}.
  \item Cumplimiento objetivo general: \texttt{docs/achievements/cumplimiento\_objetivo\_general.md}.
  \item Pipelines/simulación/análisis: \texttt{python-analysis/*.py}, \texttt{scripts/*.sh}.
\end{itemize}


\section*{Cumplimiento de Objetivos Específicos}
\noindent\textbf{Fecha:} 12 de noviembre de 2025\\

Este capítulo documenta, con formato de tesis, cómo se cumplieron los cuatro objetivos específicos del proyecto. Para cada objetivo se presenta: enunciado, enfoque, implementación, evidencia, métricas, incidencias y estado.

\subsection*{Lista de objetivos}
\begin{enumerate}
  \item Modelar la arquitectura de un prototipo de SMA capaz de orquestar y gestionar fases clave del pipeline de ML supervisado de clasificación (ingesta, preprocesamiento, preparación de datos) en entornos simulados.
  \item Implementar el prototipo de arquitectura SMA modelado, integrando al menos dos algoritmos representativos (Random Forest y SVM) para su ejecución durante la fase de entrenamiento, usando JADE como entorno multiagente.
  \item Evaluar la eficiencia computacional y la sostenibilidad energética del prototipo SMA frente a una arquitectura centralizada, durante el entrenamiento, mediante la medición de tiempo, throughput, uso de CPU y memoria, y energía normalizada (J/MB) en escenarios simulados, controlados y replicables.
  \item Analizar e interpretar los resultados para validar la eficiencia y sostenibilidad del SMA; sintetizar implicaciones técnicas, éticas y operativas; identificar desafíos, limitaciones y oportunidades de mejora.
\end{enumerate}

\subsection*{OE1 — Modelado de la arquitectura SMA}
\textbf{Enunciado.} Modelar la arquitectura de un prototipo de SMA para orquestar el pipeline de ML supervisado en entornos simulados.\\
\textbf{Enfoque.} Arquitectura JADE + Backend (FastAPI) con exposición de métricas; simulación de etapas PREPROCESS $\rightarrow$ TRAIN $\rightarrow$ EVAL con eventos JSONL.\\
\textbf{Implementación.}
\begin{itemize}
  \item Documentación de arquitectura: \texttt{docs/architecture/architecture.md}, índice \texttt{docs/architecture/README.md}.
  \item Puente con JADE: \texttt{docs/architecture/bridge\_jade.md}; protocolo de mensajes/eventos: \texttt{docs/architecture/message-protocol.md}.
  \item Scripts: \texttt{scripts/run\_jade\_experiment.sh}, \texttt{scripts/start\_stack.sh}.
\end{itemize}
\textbf{Evidencia.} Artefactos de orquestación en \texttt{data/results/\_jade\_logs/*}; métricas de progreso en \texttt{backend/api/main.py} (\texttt{/metrics}).\\
\textbf{Métricas.} \texttt{experiment\_last\_id}, \texttt{experiment\_current\_stage\_idx}, \texttt{experiment\_current\_stage\_elapsed\_seconds}.\\
\textbf{Incidencias.} Síntesis de artefactos JADE cuando faltan; normalización de documentos movidos a \texttt{docs/architecture/*}.\\
\textbf{Estado.} Cumplido.

\subsection*{OE2 — Implementación con RF y SVM (JADE)}
\textbf{Enunciado.} Implementar el prototipo SMA integrando Random Forest y SVM durante el entrenamiento con JADE.\\
\textbf{Enfoque.} Dataset reproducible y simulación de PREPROCESS/ TRAIN\_RF/ EVAL\_RF/ TRAIN\_SVM/ EVAL\_SVM, registrando eventos JSONL; JADE opera y/o se sintetizan artefactos.\\
\textbf{Implementación.}
\begin{itemize}
  \item Generador: \texttt{python-analysis/generate\_dataset.py} $\rightarrow$ \texttt{data/raw/synthetic\_classification.csv} y metadatos (SHA-256).
  \item Simulación: \texttt{python-analysis/simulate\_experiment.py}; reportes: \texttt{python-analysis/report\_run.py}.
  \item Plataforma JADE y lanzador: \texttt{agents/jade-platform/*}, \texttt{scripts/run\_jade\_experiment.sh}.
\end{itemize}
\textbf{Evidencia.} \texttt{data/results/\textless exp\_id\textgreater/} con \texttt{meta.json}, \texttt{logs/events.jsonl}, \texttt{report/report.html}.\\
\textbf{Métricas.} \texttt{train\_duration\_seconds\_last\_by\_algo}, \texttt{train\_cpu\_avg\_percent\_last\_by\_algo}, \texttt{train\_mem\_peak\_mb\_last\_by\_algo}, \texttt{train\_records\_per\_s\_last\_by\_algo}.\\
\textbf{Incidencias.} Variantes históricas \texttt{energy\_j\_*} normalizadas en \texttt{aggregate\_metrics.py}; CLASSPATH/JADE documentado en el script.\\
\textbf{Estado.} Cumplido.

\subsection*{OE3 — Evaluación de eficiencia y energía (SMA vs. centralizado)}
\textbf{Enunciado.} Evaluar eficiencia y sostenibilidad del SMA frente a una arquitectura centralizada durante el entrenamiento, con métricas (tiempo, throughput, CPU, memoria, energía J/MB) en escenarios simulados y replicables.\\
\textbf{Enfoque.} Comparación control (invocación directa) vs. tratamiento (orquestación) como proxy de arquitectura centralizada vs. SMA; dataset congelado y semillas controladas.\\
\textbf{Implementación.}
\begin{itemize}
  \item Ejecuciones reproducibles: \texttt{scripts/monitoring\_smoke.sh} y \texttt{scripts/run\_experiment.sh}.
  \item Agregación: \texttt{python-analysis/aggregate\_metrics.py} $\rightarrow$ \texttt{data/results/aggregate\_metrics.csv} con duración, CPU, memoria, throughput, \texttt{avg\_watts}, \texttt{watts\_per\_mb}, \texttt{watts\_per\_record} y Joules derivados (\\$J = W \times s\\$).
  \item Observabilidad: Prometheus + Grafana (\texttt{infra/compose/*}); dashboards versionados; dashboard mínimo en \texttt{docs/grafana/thesis\_minimal.json}.
\end{itemize}
\textbf{Evidencia.} CSV por grupo/algoritmo; paneles de performance/energía; estado del backend en \texttt{data/results/\_server\_state/*}.\\
\textbf{Métricas.} Eficiencia: \texttt{train\_duration\_seconds}, \texttt{train\_records\_per\_s\_last}, \texttt{train\_cpu\_avg\_percent\_last}, \texttt{train\_mem\_peak\_mb\_last}. Energía: \texttt{train\_avg\_watts\_last}, \texttt{watts\_per\_mb}, \texttt{watts\_per\_record}, derivación a \texttt{joules\_per\_mb}.\\
\textbf{Incidencias.} RAPL inaccesible $\rightarrow$ proxy de potencia \texttt{cpu\_power\_w\_proxy} (calibrable con \texttt{CPU\_POWER\_W}); targets Prometheus en modo host y puertos ocupados resueltos por scripts.\\
\textbf{Limitaciones.} La arquitectura “centralizada” se modela vía control vs. tratamiento; \texttt{joules\_per\_mb} se computa a partir de campos disponibles (posible añadir al CSV).\\
\textbf{Estado.} Cumplido (en alcance de prototipo).

\subsection*{OE4 — Análisis e interpretación de resultados}
\textbf{Enunciado.} Analizar e interpretar resultados para validar la eficiencia y sostenibilidad del SMA; sintetizar implicaciones técnicas, éticas y operativas; identificar desafíos/limitaciones/oportunidades.\\
\textbf{Enfoque.} Uso de CSV consolidado + reportes por ejecución + paneles de Grafana; discusión de amenazas a la validez y ética energética.\\
\textbf{Implementación.}
\begin{itemize}
  \item Consolidación: \texttt{aggregate\_metrics.py} y \texttt{data/results/aggregate\_metrics.csv}; estadísticas: \texttt{stats\_analysis.py}.
  \item Reportes por ejecución: \texttt{python-analysis/report\_run.py}.
  \item Documentación de análisis y validez: \texttt{docs/achievements/cumplimiento\_objetivo\_general.md} (apéndices), \texttt{docs/monitoring\_spec.md}, \texttt{docs/energy\_plan.md}.
\end{itemize}
\textbf{Evidencia.} Comparaciones control/tratamiento por algoritmo; discusión de incidencias reales y amenazas a la validez.\\
\textbf{Síntesis.} Técnicas (overhead vs. eficiencia; calibración de potencia), éticas (transparencia energética), operativas (scripts reproducibles, dashboards versionados).\\
\textbf{Desafíos.} Integrar RAPL por ventana en CSV; exponer \texttt{joules\_per\_mb}; extender a datasets reales y cargas intensivas.\\
\textbf{Estado.} Cumplido.

\subsection*{Conclusión}
Los cuatro objetivos específicos se encuentran cumplidos en el alcance de prototipo, con trazabilidad a código, scripts y artefactos, y con análisis cuantitativo y cualitativo suficiente para sustentar la evaluación de eficiencia y sostenibilidad.

\subsection*{Referencias cruzadas}
\begin{itemize}
  \item Arquitectura: \texttt{docs/architecture/README.md}, \texttt{architecture.md}, \texttt{bridge\_jade.md}, \texttt{message-protocol.md}.
  \item Observabilidad/energía: \texttt{docs/monitoring\_spec.md}, \texttt{docs/energy\_plan.md}, dashboards en \texttt{infra/compose/grafana/dashboards/}, \texttt{docs/grafana/README.md}.
  \item Cumplimiento objetivo general: \texttt{docs/achievements/cumplimiento\_objetivo\_general.md}.
  \item Pipelines/simulación/análisis: \texttt{python-analysis/*.py}, \texttt{scripts/*.sh}.
\end{itemize}


\section*{Cumplimiento de Objetivos Específicos}
\noindent\textbf{Fecha:} 12 de noviembre de 2025\\

Este capítulo documenta, con formato de tesis, cómo se cumplieron los cuatro objetivos específicos del proyecto. Para cada objetivo se presenta: enunciado, enfoque, implementación, evidencia, métricas, incidencias y estado.

\subsection*{Lista de objetivos}
\begin{enumerate}
  \item Modelar la arquitectura de un prototipo de SMA capaz de orquestar y gestionar fases clave del pipeline de ML supervisado de clasificación (ingesta, preprocesamiento, preparación de datos) en entornos simulados.
  \item Implementar el prototipo de arquitectura SMA modelado, integrando al menos dos algoritmos representativos (Random Forest y SVM) para su ejecución durante la fase de entrenamiento, usando JADE como entorno multiagente.
  \item Evaluar la eficiencia computacional y la sostenibilidad energética del prototipo SMA frente a una arquitectura centralizada, durante el entrenamiento, mediante la medición de tiempo, throughput, uso de CPU y memoria, y energía normalizada (J/MB) en escenarios simulados, controlados y replicables.
  \item Analizar e interpretar los resultados para validar la eficiencia y sostenibilidad del SMA; sintetizar implicaciones técnicas, éticas y operativas; identificar desafíos, limitaciones y oportunidades de mejora.
\end{enumerate}

\subsection*{OE1 — Modelado de la arquitectura SMA}
\textbf{Enunciado.} Modelar la arquitectura de un prototipo de SMA para orquestar el pipeline de ML supervisado en entornos simulados.\\
\textbf{Enfoque.} Arquitectura JADE + Backend (FastAPI) con exposición de métricas; simulación de etapas PREPROCESS $\rightarrow$ TRAIN $\rightarrow$ EVAL con eventos JSONL.\\
\textbf{Implementación.}
\begin{itemize}
  \item Documentación de arquitectura: \texttt{docs/architecture/architecture.md}, índice \texttt{docs/architecture/README.md}.
  \item Puente con JADE: \texttt{docs/architecture/bridge\_jade.md}; protocolo de mensajes/eventos: \texttt{docs/architecture/message-protocol.md}.
  \item Scripts: \texttt{scripts/run\_jade\_experiment.sh}, \texttt{scripts/start\_stack.sh}.
\end{itemize}
\textbf{Evidencia.} Artefactos de orquestación en \texttt{data/results/\_jade\_logs/*}; métricas de progreso en \texttt{backend/api/main.py} (\texttt{/metrics}).\\
\textbf{Métricas.} \texttt{experiment\_last\_id}, \texttt{experiment\_current\_stage\_idx}, \texttt{experiment\_current\_stage\_elapsed\_seconds}.\\
\textbf{Incidencias.} Síntesis de artefactos JADE cuando faltan; normalización de documentos movidos a \texttt{docs/architecture/*}.\\
\textbf{Estado.} Cumplido.

\subsection*{OE2 — Implementación con RF y SVM (JADE)}
\textbf{Enunciado.} Implementar el prototipo SMA integrando Random Forest y SVM durante el entrenamiento con JADE.\\
\textbf{Enfoque.} Dataset reproducible y simulación de PREPROCESS/ TRAIN\_RF/ EVAL\_RF/ TRAIN\_SVM/ EVAL\_SVM, registrando eventos JSONL; JADE opera y/o se sintetizan artefactos.\\
\textbf{Implementación.}
\begin{itemize}
  \item Generador: \texttt{python-analysis/generate\_dataset.py} $\rightarrow$ \texttt{data/raw/synthetic\_classification.csv} y metadatos (SHA-256).
  \item Simulación: \texttt{python-analysis/simulate\_experiment.py}; reportes: \texttt{python-analysis/report\_run.py}.
  \item Plataforma JADE y lanzador: \texttt{agents/jade-platform/*}, \texttt{scripts/run\_jade\_experiment.sh}.
\end{itemize}
\textbf{Evidencia.} \texttt{data/results/\textless exp\_id\textgreater/} con \texttt{meta.json}, \texttt{logs/events.jsonl}, \texttt{report/report.html}.\\
\textbf{Métricas.} \texttt{train\_duration\_seconds\_last\_by\_algo}, \texttt{train\_cpu\_avg\_percent\_last\_by\_algo}, \texttt{train\_mem\_peak\_mb\_last\_by\_algo}, \texttt{train\_records\_per\_s\_last\_by\_algo}.\\
\textbf{Incidencias.} Variantes históricas \texttt{energy\_j\_*} normalizadas en \texttt{aggregate\_metrics.py}; CLASSPATH/JADE documentado en el script.\\
\textbf{Estado.} Cumplido.

\subsection*{OE3 — Evaluación de eficiencia y energía (SMA vs. centralizado)}
\textbf{Enunciado.} Evaluar eficiencia y sostenibilidad del SMA frente a una arquitectura centralizada durante el entrenamiento, con métricas (tiempo, throughput, CPU, memoria, energía J/MB) en escenarios simulados y replicables.\\
\textbf{Enfoque.} Comparación control (invocación directa) vs. tratamiento (orquestación) como proxy de arquitectura centralizada vs. SMA; dataset congelado y semillas controladas.\\
\textbf{Implementación.}
\begin{itemize}
  \item Ejecuciones reproducibles: \texttt{scripts/monitoring\_smoke.sh} y \texttt{scripts/run\_experiment.sh}.
  \item Agregación: \texttt{python-analysis/aggregate\_metrics.py} $\rightarrow$ \texttt{data/results/aggregate\_metrics.csv} con duración, CPU, memoria, throughput, \texttt{avg\_watts}, \texttt{watts\_per\_mb}, \texttt{watts\_per\_record} y Joules derivados (\\$J = W \times s\\$).
  \item Observabilidad: Prometheus + Grafana (\texttt{infra/compose/*}); dashboards versionados; dashboard mínimo en \texttt{docs/grafana/thesis\_minimal.json}.
\end{itemize}
\textbf{Evidencia.} CSV por grupo/algoritmo; paneles de performance/energía; estado del backend en \texttt{data/results/\_server\_state/*}.\\
\textbf{Métricas.} Eficiencia: \texttt{train\_duration\_seconds}, \texttt{train\_records\_per\_s\_last}, \texttt{train\_cpu\_avg\_percent\_last}, \texttt{train\_mem\_peak\_mb\_last}. Energía: \texttt{train\_avg\_watts\_last}, \texttt{watts\_per\_mb}, \texttt{watts\_per\_record}, derivación a \texttt{joules\_per\_mb}.\\
\textbf{Incidencias.} RAPL inaccesible $\rightarrow$ proxy de potencia \texttt{cpu\_power\_w\_proxy} (calibrable con \texttt{CPU\_POWER\_W}); targets Prometheus en modo host y puertos ocupados resueltos por scripts.\\
\textbf{Limitaciones.} La arquitectura “centralizada” se modela vía control vs. tratamiento; \texttt{joules\_per\_mb} se computa a partir de campos disponibles (posible añadir al CSV).\\
\textbf{Estado.} Cumplido (en alcance de prototipo).

\subsection*{OE4 — Análisis e interpretación de resultados}
\textbf{Enunciado.} Analizar e interpretar resultados para validar la eficiencia y sostenibilidad del SMA; sintetizar implicaciones técnicas, éticas y operativas; identificar desafíos/limitaciones/oportunidades.\\
\textbf{Enfoque.} Uso de CSV consolidado + reportes por ejecución + paneles de Grafana; discusión de amenazas a la validez y ética energética.\\
\textbf{Implementación.}
\begin{itemize}
  \item Consolidación: \texttt{aggregate\_metrics.py} y \texttt{data/results/aggregate\_metrics.csv}; estadísticas: \texttt{stats\_analysis.py}.
  \item Reportes por ejecución: \texttt{python-analysis/report\_run.py}.
  \item Documentación de análisis y validez: \texttt{docs/achievements/cumplimiento\_objetivo\_general.md} (apéndices), \texttt{docs/monitoring\_spec.md}, \texttt{docs/energy\_plan.md}.
\end{itemize}
\textbf{Evidencia.} Comparaciones control/tratamiento por algoritmo; discusión de incidencias reales y amenazas a la validez.\\
\textbf{Síntesis.} Técnicas (overhead vs. eficiencia; calibración de potencia), éticas (transparencia energética), operativas (scripts reproducibles, dashboards versionados).\\
\textbf{Desafíos.} Integrar RAPL por ventana en CSV; exponer \texttt{joules\_per\_mb}; extender a datasets reales y cargas intensivas.\\
\textbf{Estado.} Cumplido.

\subsection*{Conclusión}
Los cuatro objetivos específicos se encuentran cumplidos en el alcance de prototipo, con trazabilidad a código, scripts y artefactos, y con análisis cuantitativo y cualitativo suficiente para sustentar la evaluación de eficiencia y sostenibilidad.

\subsection*{Referencias cruzadas}
\begin{itemize}
  \item Arquitectura: \texttt{docs/architecture/README.md}, \texttt{architecture.md}, \texttt{bridge\_jade.md}, \texttt{message-protocol.md}.
  \item Observabilidad/energía: \texttt{docs/monitoring\_spec.md}, \texttt{docs/energy\_plan.md}, dashboards en \texttt{infra/compose/grafana/dashboards/}, \texttt{docs/grafana/README.md}.
  \item Cumplimiento objetivo general: \texttt{docs/achievements/cumplimiento\_objetivo\_general.md}.
  \item Pipelines/simulación/análisis: \texttt{python-analysis/*.py}, \texttt{scripts/*.sh}.
\end{itemize}
